\chapter{Penis Care}
\index{Penis Care}

\medskip
\noindent\textit{This chapter is for general education. Persistent pain, discharge, sores, or changes in urination need clinical assessment.}

\section*{Routine care}
Wash the penis gently with warm water and dry it. If the foreskin retracts comfortably, clean underneath and replace it over the glans afterwards. Do not force a child’s foreskin back. Avoid perfumed soap, deodorant, harsh scrubbing, and products that irritate the skin.

\section*{Sexual and urinary health}
Use condoms or other barriers according to sexual-health advice, and seek STI testing after a relevant exposure or when symptoms occur. Do not share medicines. Painful urination, discharge, sores, new lumps, blood in urine or semen, or changes in erections should be assessed rather than assumed to be an infection.

\section*{Foreskin problems}
A tight foreskin may be normal in young children but should not be forced. In adults or children, recurrent swelling, bleeding, discharge, pain, or difficulty urinating may indicate inflammation, infection, scarring, or another condition.

\section*{When to Seek Medical Care}
Seek emergency help if the foreskin is stuck behind the glans and the penis is swelling or painful, if you cannot pass urine, or for severe injury, uncontrolled bleeding, or rapidly increasing swelling. Arrange medical review for persistent irritation, discharge, sores, pain, or a new lump.

\section*{Sources and further reading}
\begin{itemize}
\item NHS. \textit{Balanitis}. \url{https://www.nhs.uk/conditions/balanitis/}
\item NHS. \textit{Tight foreskin (phimosis)}. \url{https://www.nhs.uk/conditions/phimosis/}
\end{itemize}
